\chapter{Appendix: Algebra, Functions, and Proof}

Calculus does not begin from nothing. Every limit you compute, every derivative
you simplify, every integral you wrestle into a closed form leans on a quieter
layer of mathematics underneath: the algebra of rearranging expressions, the
laws of exponents, the handling of inequalities and absolute values, a few
durable facts about functions, and the small but crucial habit of \emph{proving}
that something is true rather than merely believing it. This appendix gathers
that layer in one place. It is the toolbox you reach into without thinking once
you are fluent --- and the place to come back to when a step in a derivation
suddenly feels like a magic trick.

The spirit here is the same as the rest of the book. We will not just list rules.
We will ask \emph{why} each rule holds, draw the picture that makes it obvious,
and point out the traps where good students reliably slip. By the end you should
feel that none of this is arbitrary: the exponent laws are forced on us once we
decide what $a^n$ should mean; the rules for inequalities follow from one honest
fact about the number line; and a proof is not a ritual but an explanation that
happens to be airtight.

\begin{focusbox}
Three things to carry away from this appendix:
\begin{enumerate}
  \item Algebraic identities (factoring, exponent laws, the binomial pattern)
        are not memorized magic --- each is forced by a single defining idea, and
        once you see the idea you can reconstruct the rule.
  \item Inequalities and absolute values are the native language of calculus
        (limits, error bounds, convergence). The one rule everyone forgets ---
        multiplying by a negative flips the sign --- has a clean reason.
  \item A proof is an explanation with no gaps. Direct proof, proof by
        contradiction, and induction are three reliable shapes that explanation
        can take.
\end{enumerate}
\end{focusbox}

\section{Numbers, and the Ground We Stand On}
\label{ch40:numbers}

Before algebra there are the numbers algebra rearranges. We use a nested tower of
number systems, each built to fix a shortcoming of the last:
\[
  \N \subset \Z \subset \Q \subset \R \subset \C.
\]
The naturals $\N = \set{1,2,3,\dots}$ count things. The integers
$\Z = \set{\dots,-2,-1,0,1,2,\dots}$ add the ability to subtract freely. The
rationals $\Q$ --- ratios $p/q$ of integers with $q \neq 0$ --- add the ability to
divide. The reals $\R$ fill in the gaps between rationals so that the number line
has no holes (a number like $\sqrt 2$ or $\pi$ lives here but not in $\Q$). The
complex numbers $\C$ add a square root of $-1$. Calculus lives almost entirely in
$\R$, and the ``no holes'' property of $\R$ --- its \emph{completeness} --- is
what makes limits behave; that story has its own appendix.

\begin{intuitionbox}
Picture the real line as an infinitely long, perfectly smooth ruler. Every point
on it is a real number, and every real number sits at exactly one point. The
rationals are dense --- between any two of them there is another --- yet they leave
gaps you cannot see but can prove are there. The reals are what you get when you
plug every gap. Most of calculus is the art of saying precise things about points
sitting infinitely close together on this ruler.
\end{intuitionbox}

The arithmetic of $\R$ obeys a short list of field axioms that we will use without
ceremony: addition and multiplication are commutative and associative, $0$ and $1$
are identities, every number has an additive inverse and every nonzero number a
multiplicative inverse, and multiplication distributes over addition:
\[
  a(b+c) = ab + ac.
\]
That last law, the \emph{distributive law}, is the single most overworked fact in
all of algebra. Expanding brackets, factoring, multiplying polynomials, the
binomial theorem --- all of it is the distributive law applied again and again.

\begin{happeningbox}
When you ``expand $3(x+2)$ to get $3x+6$'' and when you ``factor $3x+6$ back to
$3(x+2)$'' you are doing the exact same equation read in two directions:
$a(b+c) = ab+ac$. Expanding goes left to right; factoring goes right to left.
This is why factoring feels harder --- running a rule \emph{backwards} means
searching for the inputs that would have produced what you see, and search is
always harder than evaluation.
\end{happeningbox}

\section{Algebraic Manipulation: Expanding and Factoring}
\label{ch40:factoring}

\subsection{Multiplying out}

To multiply two sums, distribute every term of the first across every term of the
second. For two binomials this is the familiar pattern (sometimes taught as
``FOIL,'' though the name hides the principle):
\[
  (a+b)(c+d) = ac + ad + bc + bd.
\]
Each term on the right is one term from the first factor times one term from the
second; we collect every such pairing. There is nothing special about two terms.
$(a+b)(c+d+e)$ produces six products, one for each (left term, right term) pair.

The identities worth knowing cold are the ones that appear constantly in
calculus, especially when rationalizing or simplifying difference quotients:
\begin{align}
  (a+b)^2 &= a^2 + 2ab + b^2, \label{ch40:square-sum}\\
  (a-b)^2 &= a^2 - 2ab + b^2, \\
  (a+b)(a-b) &= a^2 - b^2, \label{ch40:diff-squares}\\
  (a+b)^3 &= a^3 + 3a^2 b + 3ab^2 + b^3.
\end{align}
The difference of squares~\eqref{ch40:diff-squares} is the workhorse: it is the
engine behind the trick of multiplying by a conjugate, which you meet the very
first time you compute a limit like $\lim_{x\to 0}(\sqrt{x+1}-1)/x$.

\begin{whybox}
Why is $(a+b)^2 = a^2 + 2ab + b^2$ and \emph{not} $a^2 + b^2$? Because squaring a
sum is not the same as summing the squares --- the cross terms are real. Run the
distributive law honestly:
\[
  (a+b)^2 = (a+b)(a+b) = a\cdot a + a \cdot b + b \cdot a + b \cdot b
          = a^2 + 2ab + b^2.
\]
The middle $2ab$ is the two ``mixed'' pairings $ab$ and $ba$. Forgetting it is
probably the single most common algebra error in all of calculus, and the
picture below shows exactly what it is.
\end{whybox}

The cross term is not abstract. It is literally the two rectangles you get when
you enlarge a square.

\begin{figure}[H]
\centering
\begin{tikzpicture}[scale=1.1]
  % big square split into 4 regions
  \draw[fill=sourceblue!18] (0,0) rectangle (3,3);
  \draw[fill=quotientorange!25] (3,0) rectangle (4.2,3);
  \draw[fill=quotientorange!25] (0,3) rectangle (3,4.2);
  \draw[fill=targetgreen!22] (3,3) rectangle (4.2,4.2);
  \draw[thick] (0,0) rectangle (4.2,4.2);
  \draw[thick] (3,0) -- (3,4.2);
  \draw[thick] (0,3) -- (4.2,3);
  % labels of regions
  \node at (1.5,1.5) {$a^2$};
  \node at (3.6,1.5) {$ab$};
  \node at (1.5,3.6) {$ab$};
  \node at (3.6,3.6) {$b^2$};
  % side labels
  \node[below] at (1.5,-0.05) {$a$};
  \node[below] at (3.6,-0.05) {$b$};
  \node[left]  at (-0.05,1.5) {$a$};
  \node[left]  at (-0.05,3.6) {$b$};
\end{tikzpicture}
\caption{The square on side $a+b$ has area $(a+b)^2$. Cut it into four pieces:
the $a$-square, the $b$-square, and \emph{two} $a\times b$ rectangles. That is why
$(a+b)^2 = a^2 + 2ab + b^2$ --- the $2ab$ is the two orange rectangles, the part
people forget.}
\label{ch40:fig-square}
\end{figure}

\subsection{Factoring: running the machine backwards}

Factoring rewrites a sum as a product. It matters because a product is zero
exactly when one of its factors is zero (the \emph{zero-product property}), which
is how we solve equations, and because a factored expression often lets a
troublesome common factor cancel --- the heart of computing limits of the
indeterminate form $0/0$.

The basic moves, in rough order of how often you need them:

\textbf{(1) Pull out a common factor.}
\[
  6x^3 - 9x^2 = 3x^2(2x - 3).
\]
Look for the largest factor shared by every term and lift it out front.

\textbf{(2) Difference of squares.}
\[
  x^2 - 25 = (x-5)(x+5), \qquad x^4 - 1 = (x^2-1)(x^2+1) = (x-1)(x+1)(x^2+1).
\]

\textbf{(3) Quadratic trinomials.} To factor $x^2 + bx + c$, find two numbers that
\emph{multiply to $c$ and add to $b$}. For $x^2 + 7x + 12$ we want a pair with
product $12$ and sum $7$: that is $3$ and $4$, so
\[
  x^2 + 7x + 12 = (x+3)(x+4).
\]

\textbf{(4) Sum and difference of cubes.}
\[
  a^3 - b^3 = (a-b)(a^2 + ab + b^2), \qquad
  a^3 + b^3 = (a+b)(a^2 - ab + b^2).
\]

\begin{trapbox}
There is no real factoring of a \emph{sum} of squares: $a^2 + b^2$ does not factor
over the reals (you need complex numbers). Students who just factored $a^2 - b^2$
often write $a^2+b^2 = (a+b)(a-b)$ out of momentum --- but expanding that gives
$a^2 - b^2$, the wrong sign. The difference factors; the sum does not.
\end{trapbox}

\begin{example}[Factoring to kill a $0/0$ limit]
\label{ch40:ex-factor-limit}
Consider $\dfrac{x^2 - 9}{x - 3}$ as $x \to 3$. Plugging in $x=3$ gives $0/0$,
which says nothing yet. Factor the top as a difference of squares:
\[
  \frac{x^2 - 9}{x-3} = \frac{(x-3)(x+3)}{x-3} = x + 3 \quad (x \neq 3).
\]
The offending factor $x-3$ cancels, and the simplified expression $x+3$ tends to
$6$ as $x \to 3$. The whole drama of the limit was a single hidden common factor;
factoring exposed and removed it.
\end{example}

\begin{walkbox}
Try factoring $2x^2 + x - 6$. (Hint: look for two numbers multiplying to
$2\cdot(-6) = -12$ and adding to $1$; they are $4$ and $-3$. Split the middle term:
$2x^2 + 4x - 3x - 6 = 2x(x+2) - 3(x+2) = (2x-3)(x+2)$.) The technique of
``splitting the middle term'' turns a hard trinomial into two easy common-factor
pulls.
\end{walkbox}

\subsection{The quadratic formula, and completing the square}

Not every quadratic factors with nice integers. For \emph{any} quadratic
$ax^2 + bx + c = 0$ with $a \neq 0$, the solutions are
\[
  x = \frac{-b \pm \sqrt{b^2 - 4ac}}{2a}.
\]
This is not handed down from the sky --- it comes from \emph{completing the
square}, a technique worth seeing in full because the same maneuver reappears in
integration (trig substitution) and in finding vertices of parabolas.

\begin{whybox}
Start from $ax^2 + bx + c = 0$. Divide by $a$ and move the constant:
\[
  x^2 + \tfrac{b}{a}x = -\tfrac{c}{a}.
\]
The left side is \emph{almost} a perfect square. The square $(x + \tfrac{b}{2a})^2$
expands to $x^2 + \tfrac{b}{a}x + \tfrac{b^2}{4a^2}$, so add $\tfrac{b^2}{4a^2}$ to
both sides to complete it:
\[
  \Bigl(x + \tfrac{b}{2a}\Bigr)^2 = \frac{b^2}{4a^2} - \frac{c}{a}
                                  = \frac{b^2 - 4ac}{4a^2}.
\]
Now take square roots (this is where the $\pm$ is born) and solve for $x$:
\[
  x + \frac{b}{2a} = \pm\frac{\sqrt{b^2-4ac}}{2a}
  \;\Longrightarrow\;
  x = \frac{-b \pm \sqrt{b^2-4ac}}{2a}.
\]
The quadratic formula is just completing the square done once, symbolically, so
you never have to do it again.
\end{whybox}

\begin{historybox}
The technique of completing the square is ancient. Babylonian scribes were solving
quadratics by an equivalent geometric ``cut-and-paste'' procedure nearly four
thousand years ago. Around 825 CE the Persian mathematician al-Khwarizmi wrote
\emph{al-Kitab al-mukhtasar fi hisab al-jabr wa'l-muqabala} --- ``the compendious
book on calculation by restoration and balancing.'' The word \emph{al-jabr}
(``restoration,'' meaning moving a subtracted term to the other side) is the root
of our word \emph{algebra}. He literally completed squares by drawing them, the
same picture as Figure~\ref{ch40:fig-square}.
\end{historybox}

The quantity $b^2 - 4ac$ under the root is the \emph{discriminant}. If it is
positive there are two real roots; if zero, one repeated root; if negative, no real
roots (the parabola misses the $x$-axis) and the roots are complex conjugates.

\section{Exponents and Their Laws}
\label{ch40:exponents}

\subsection{Where the laws come from}

For a positive integer $n$, $a^n$ means ``multiply $a$ by itself $n$ times.'' From
this single definition every exponent law follows by counting factors. Take
$a^m \cdot a^n$: that is $m$ copies of $a$ times $n$ copies of $a$, which is $m+n$
copies. Hence
\[
  a^m \cdot a^n = a^{m+n}.
\]
There is nothing to memorize; you are just counting how many $a$'s stand in the
line. The same counting gives the rest:
\begin{align}
  a^m \cdot a^n &= a^{m+n}, & \frac{a^m}{a^n} &= a^{m-n}, \\
  (a^m)^n &= a^{mn}, & (ab)^n &= a^n b^n, \\
  \Bigl(\frac{a}{b}\Bigr)^n &= \frac{a^n}{b^n}. & &
\end{align}

\begin{happeningbox}
Every exponent law is really the \emph{single} fact $a^{m+n}=a^m a^n$ wearing
different clothes, because exponentiation turns addition (in the exponent) into
multiplication (of the values). $(a^m)^n = a^{mn}$ is repeated addition of
exponents: $n$ copies of $m$. This addition-becomes-multiplication structure is
exactly what makes the exponential function $e^x$ so central in calculus --- it is
the function that converts the additive world of inputs into the multiplicative
world of outputs, and it is the only one (up to scaling the input) that is its own
derivative.
\end{happeningbox}

\subsection{Extending to zero, negatives, and fractions}

The definition ``multiply $a$ by itself $n$ times'' only makes sense for positive
integer $n$. So what could $a^0$, $a^{-3}$, or $a^{1/2}$ possibly mean? The
guiding principle is beautiful: \emph{we extend the definition in the only way that
keeps the law $a^m a^n = a^{m+n}$ true.} The laws are not re-proved for the new
cases; they are the very thing we insist on preserving, and that insistence forces
a unique answer.

\textbf{Zero exponent.} We need $a^0 \cdot a^n = a^{0+n} = a^n$. Dividing by $a^n$
(for $a \neq 0$) forces
\[
  a^0 = 1.
\]

\textbf{Negative exponent.} We need $a^{-n} \cdot a^n = a^{0} = 1$, so $a^{-n}$
must be the reciprocal:
\[
  a^{-n} = \frac{1}{a^n}.
\]

\textbf{Fractional exponent.} We need $\bigl(a^{1/n}\bigr)^n = a^{(1/n)\cdot n} =
a^1 = a$, so $a^{1/n}$ must be the number whose $n$th power is $a$ --- the $n$th
root:
\[
  a^{1/n} = \sqrt[n]{a}, \qquad a^{m/n} = \bigl(\sqrt[n]{a}\bigr)^m = \sqrt[n]{a^m}.
\]

\begin{whybox}
This is one of the most important moves in all of mathematics, so let it land.
The expression $a^{1/2}$ has \emph{no} prior meaning --- you cannot multiply $a$
by itself ``half a time.'' We \emph{define} it. But we do not define it by whim:
we choose the one value that lets the exponent law survive. That demand pins it
down uniquely to $\sqrt a$. Mathematics grows by exactly this method: take a rule
that holds in a small world, then extend the world in whatever way keeps the rule
intact. The same logic later defines $a^x$ for irrational $x$, $0! = 1$, and
$e^{i\theta} = \cos\theta + i\sin\theta$.
\end{whybox}

\begin{figure}[H]
\centering
\begin{tikzpicture}
\begin{axis}[
    width=11.5cm, height=6.4cm,
    xlabel={$x$}, ylabel={$2^x$},
    domain=-3:3, samples=120,
    xmin=-3.2, xmax=3.2, ymin=-0.4, ymax=8.4,
    axis lines=middle,
    grid=major, grid style={gray!18},
    legend pos=north west,
]
  \addplot[curveblue, very thick] {2^x};
  \addlegendentry{$y = 2^x$}
  % mark special integer points
  \addplot[only marks, mark=*, mark size=1.6pt, tangentred]
     coordinates {(-2,0.25) (-1,0.5) (0,1) (1,2) (2,4) (3,8)};
  \node[anchor=south west, font=\small] at (axis cs:0,1) {$2^0=1$};
  \node[anchor=west, font=\small] at (axis cs:-2,0.25) {$2^{-2}=\tfrac14$};
\end{axis}
\end{tikzpicture}
\caption{The exponential $2^x$ knits the integer powers (red dots) into a smooth
curve. Negative exponents give reciprocals ($2^{-2}=\tfrac14$), the zero exponent
gives $1$, and fractional exponents fill in everything between. The values for
non-integer $x$ are precisely the ones that keep $2^{m+n}=2^m2^n$ holding for all
real $m,n$.}
\label{ch40:fig-exp}
\end{figure}

\begin{trapbox}
The exponent laws govern \emph{products with a shared base}, not sums and not
mixed bases. All of the following are wrong:
\[
  a^m + a^n \neq a^{m+n}, \qquad
  (a+b)^n \neq a^n + b^n, \qquad
  \sqrt{a^2+b^2} \neq a + b.
\]
The last is the same error as the missing cross term in $(a+b)^2$. There is no law
that simplifies $a^m + a^n$ in general --- the laws live on the multiplicative side
only. Whenever you are tempted to ``distribute'' an exponent across a sum, stop:
exponents distribute over products, never over sums.
\end{trapbox}

\section{Inequalities}
\label{ch40:inequalities}

Equations ask \emph{where} two things are equal; inequalities ask \emph{which is
bigger}. Calculus runs on inequalities far more than on equations: the definition
of a limit is an inequality ($\abs{f(x)-L}<\eps$), comparison tests for series are
inequalities, error bounds are inequalities. So fluency here pays off everywhere.

The order relation $<$ on $\R$ obeys a handful of rules. Most are unsurprising, but
one is the famous trap.

\begin{focusbox}
\textbf{Rules for inequalities.} For real numbers,
\begin{itemize}
  \item \textbf{Add anything:} if $a < b$ then $a + c < b + c$ (for any $c$).
  \item \textbf{Multiply by a positive:} if $a<b$ and $c>0$ then $ac < bc$.
  \item \textbf{Multiply by a negative --- FLIP:} if $a<b$ and $c<0$ then
        $ac > bc$.
  \item \textbf{Transitivity:} if $a<b$ and $b<c$ then $a<c$.
  \item \textbf{Reciprocal of positives flips:} if $0<a<b$ then
        $\tfrac1a > \tfrac1b$.
\end{itemize}
\end{focusbox}

\begin{whybox}
Why does multiplying by a negative \emph{flip} the inequality? Picture the number
line. The statement $a<b$ means $a$ sits to the left of $b$. Multiplying every
number by $-1$ is a \emph{reflection of the line through $0$}: left and right swap.
So whatever was on the left ($a$) lands on the right, and the order reverses. More
algebraically: $a<b$ means $b-a>0$, i.e.\ $b-a$ is positive. Multiply a positive
number by a negative $c$ and you get a negative: $c(b-a)<0$, that is $cb-ca<0$,
i.e.\ $cb<ca$, i.e.\ $ca>cb$. The flip is forced.
\end{whybox}

\begin{figure}[H]
\centering
\begin{tikzpicture}[scale=1.0]
  % original line
  \draw[->,thick] (-4.6,0) -- (4.6,0);
  \foreach \x in {-4,...,4} \draw (\x,0.08)--(\x,-0.08);
  \fill[sourceblue] (1,0) circle (2.6pt) node[above=3pt]{$a=1$};
  \fill[targetgreen] (3,0) circle (2.6pt) node[above=3pt]{$b=3$};
  \node at (-5.4,0) {original};
  % reflected line below
  \draw[->,thick] (-4.6,-1.8) -- (4.6,-1.8);
  \foreach \x in {-4,...,4} \draw (\x,-1.72)--(\x,-1.88);
  \fill[sourceblue] (-1,-1.8) circle (2.6pt) node[below=3pt]{$-a=-1$};
  \fill[targetgreen] (-3,-1.8) circle (2.6pt) node[below=3pt]{$-b=-3$};
  \node at (-5.4,-1.8) {$\times(-1)$};
  % reflection arrows
  \draw[->,trapred,dashed] (1,-0.18) to[bend right=25] (-1,-1.62);
  \draw[->,trapred,dashed] (3,-0.18) to[bend right=25] (-3,-1.62);
\end{tikzpicture}
\caption{Multiplying by $-1$ reflects the line through $0$. The blue point $a=1$
was left of the green $b=3$; after reflection $-a=-1$ is to the \emph{right} of
$-b=-3$. The order has flipped: $1<3$ but $-1>-3$.}
\label{ch40:fig-flip}
\end{figure}

\begin{trapbox}
The flip rule is the most common source of wrong inequality answers. When you
divide both sides of $-2x < 6$ by $-2$, you must reverse: $x > -3$, \emph{not}
$x<-3$. A safe habit: never multiply or divide an inequality by a quantity whose
sign you do not know. If you must divide by something like $x-1$, split into cases
$x>1$ and $x<1$ and handle each separately.
\end{trapbox}

\begin{example}[Solving a rational inequality the safe way]
\label{ch40:ex-rational-ineq}
Solve $\dfrac{x+1}{x-2} \ge 0$. Do \emph{not} cross-multiply --- you do not know
the sign of $x-2$. Instead find where each factor is zero or undefined: the
numerator vanishes at $x=-1$, the denominator at $x=2$ (where the expression is
undefined). These split the line into three regions. Test one point in each:
\[
  \begin{array}{c|ccc}
    \text{region} & x<-1 & -1<x<2 & x>2 \\\hline
    \text{sign of } \tfrac{x+1}{x-2} & (-)/(-)=+ & (+)/(-)=- & (+)/(+)=+
  \end{array}
\]
So the expression is $\ge 0$ on $x \le -1$ and $x > 2$ (including $x=-1$ where it
equals $0$, excluding $x=2$ where it is undefined). The solution is
$(-\infty,-1] \cup (2,\infty)$.
\end{example}

The whole method is a \emph{sign chart}: mark the points where the expression is
zero or undefined, then record the sign in each gap. Figure~\ref{ch40:fig-signchart}
shows it at a glance.

\begin{figure}[H]
\centering
\begin{tikzpicture}[scale=1.0]
  \draw[->,thick] (-5,0)--(5,0) node[right]{$x$};
  % critical points
  \fill[targetgreen] (-2,0) circle (2.6pt) node[below=4pt]{$x=-1$};
  \draw[fill=white,trapred,thick] (2,0) circle (2.6pt) node[below=4pt]{$x=2$};
  % zero / undefined annotations
  \node[above=3pt, targetgreen, font=\small] at (-2,0) {zero};
  \node[above=3pt, trapred, font=\small] at (2,0) {undefined};
  % signs in the three regions
  \node[sourceblue, font=\large] at (-3.6,0.45) {$+$};
  \node[trapred,    font=\large] at (0,0.45)    {$-$};
  \node[sourceblue, font=\large] at (3.6,0.45)  {$+$};
  % shaded solution set
  \draw[sourceblue, line width=2.6pt] (-5,-0.55)--(-2,-0.55);
  \fill[sourceblue] (-2,-0.55) circle (2.6pt);
  \draw[sourceblue, line width=2.6pt] (2,-0.55)--(5,-0.55);
  \draw[fill=white,sourceblue,thick] (2,-0.55) circle (2.6pt);
  \node[sourceblue, font=\small, anchor=west] at (-5,-1.05)
       {solution of $\tfrac{x+1}{x-2}\ge 0$: $(-\infty,-1]\cup(2,\infty)$};
\end{tikzpicture}
\caption{Sign chart for $\tfrac{x+1}{x-2}$. The two critical points --- the zero at
$x=-1$ (filled, included) and the undefined point at $x=2$ (open, excluded) --- cut
the line into three regions. One test point per region fixes the sign there. The
heavy blue segments below the axis are the solution set, where the fraction is
$\ge 0$.}
\label{ch40:fig-signchart}
\end{figure}

\section{Absolute Value: Distance Made Algebraic}
\label{ch40:absolute-value}

The absolute value $\abs{x}$ is defined by cases,
\[
  \abs{x} = \begin{cases} x, & x \ge 0,\\ -x, & x < 0,\end{cases}
\]
but the case split hides the real idea. The honest reading is geometric:

\begin{intuitionbox}
$\abs{x}$ is the \emph{distance from $x$ to $0$} on the number line, and more
generally $\abs{x-a}$ is the \emph{distance between $x$ and $a$}. The two cases
just say ``distance is never negative; if the number is already to the right of
$0$ its distance is itself, and if it is to the left we negate to make it
positive.'' Read every absolute value as a distance and most manipulations become
obvious.
\end{intuitionbox}

This distance reading immediately explains how absolute-value inequalities unfold:
\[
  \abs{x-a} < r \quad\Longleftrightarrow\quad a - r < x < a + r,
\]
``$x$ is within $r$ of $a$,'' an \emph{interval} centered at $a$ with radius $r$.
And the opposite,
\[
  \abs{x-a} > r \quad\Longleftrightarrow\quad x < a-r \;\text{ or }\; x > a+r,
\]
``$x$ is farther than $r$ from $a$,'' the \emph{outside} of that interval.

\begin{figure}[H]
\centering
\begin{tikzpicture}[scale=1.0]
  \draw[->,thick] (-1,0)--(9,0) node[right]{$x$};
  \foreach \x/\lab in {1/{a-r}, 4/{a}, 7/{a+r}}
     \draw (\x,0.09)--(\x,-0.09) node[below=2pt]{$\lab$};
  % shaded interval
  \draw[sourceblue, line width=3pt] (1,0)--(7,0);
  \fill[white,draw=sourceblue,thick] (1,0) circle (3pt);
  \fill[white,draw=sourceblue,thick] (7,0) circle (3pt);
  \fill[targetgreen] (4,0) circle (2.4pt);
  \node[above, sourceblue] at (4,0.35) {$\abs{x-a}<r$ : the open interval $(a-r,\,a+r)$};
  % radius braces
  \draw[decorate,decoration={brace,amplitude=5pt},quotientorange]
       (4,-0.7)--(1,-0.7) node[midway,below=4pt]{$r$};
  \draw[decorate,decoration={brace,amplitude=5pt},quotientorange]
       (7,-0.7)--(4,-0.7) node[midway,below=4pt]{$r$};
\end{tikzpicture}
\caption{$\abs{x-a}<r$ reads as ``$x$ is within distance $r$ of $a$,'' so the
solution set is the open interval $(a-r,a+r)$ --- a fixed center $a$ with a radius
$r$ in each direction. This exact picture is the geometry behind the
$\eps$-$\delta$ definition of a limit.}
\label{ch40:fig-absval}
\end{figure}

\begin{connectionbox}
This is not a side topic --- it is the whole shape of a limit. The statement
$\lim_{x\to a} f(x) = L$ unpacks to: for every tolerance $\eps>0$ there is a
$\delta>0$ so that $\abs{x-a}<\delta$ forces $\abs{f(x)-L}<\eps$. Both halves are
absolute-value inequalities, i.e.\ ``within distance'' statements: \emph{whenever
$x$ is within $\delta$ of $a$, $f(x)$ is within $\eps$ of $L$.}
Figure~\ref{ch40:fig-absval} is, quite literally, the inner machinery of the
central definition of calculus.
\end{connectionbox}

Two algebraic facts about absolute value earn their keep constantly.

\begin{focusbox}
\textbf{Useful identities.}
\[
  \abs{xy} = \abs{x}\,\abs{y}, \qquad \abs{x}^2 = x^2, \qquad \sqrt{x^2}=\abs{x}.
\]
\textbf{The triangle inequality.} For all real (indeed all complex) $x,y$,
\[
  \abs{x+y} \le \abs{x} + \abs{y}.
\]
\end{focusbox}

The triangle inequality is arguably the most-used inequality in all of analysis.
Its name comes from the picture in the plane: the straight side of a triangle is no
longer than the sum of the other two sides. On the line it says a detour through
the origin is never shorter than going direct.

\begin{whybox}
Here is a clean proof of the triangle inequality. For any real $t$, $-\abs t \le t
\le \abs t$ (a number lies between $\pm$ its size). Apply this to $x$ and to $y$
and add:
\[
  -\abs x \le x \le \abs x, \qquad -\abs y \le y \le \abs y
\]
add the left and right ends:
\[
  -(\abs x + \abs y) \le x + y \le \abs x + \abs y.
\]
A statement of the form $-c \le s \le c$ (with $c \ge 0$) is exactly
$\abs s \le c$. With $s = x+y$ and $c = \abs x + \abs y$ this reads
$\abs{x+y} \le \abs x + \abs y$. Done. The whole proof is the single observation
that every number is sandwiched by $\pm$ its absolute value.
\end{whybox}

\section{Functions: The Facts Calculus Leans On}
\label{ch40:functions}

Functions get a full chapter early in the book; here we collect, for reference, the
handful of facts the later chapters use without re-deriving.

\textbf{A function} $f\colon D \to \R$ assigns to each input $x$ in its domain $D$
\emph{exactly one} output $f(x)$. The ``exactly one'' is the defining demand --- it
is what the \emph{vertical line test} checks on a graph: a curve is the graph of a
function precisely when no vertical line meets it twice.

\textbf{Composition.} $(g\circ f)(x) = g(f(x))$: do $f$ first, then feed the result
to $g$. Order matters; $g\circ f$ and $f\circ g$ are usually different functions.
The chain rule is the derivative of a composition, so composition is not a
curiosity --- it is half of differential calculus.

\textbf{Inverse functions.} $f^{-1}$ undoes $f$: $f^{-1}(f(x)) = x$ and
$f(f^{-1}(y)) = y$. An inverse exists exactly when $f$ is \emph{one-to-one} (each
output comes from only one input), which the \emph{horizontal line test} checks.
The graph of $f^{-1}$ is the graph of $f$ reflected across the line $y=x$, because
inverting swaps the roles of input and output, i.e.\ swaps the axes.

\begin{figure}[H]
\centering
\begin{tikzpicture}
\begin{axis}[
    width=8.4cm, height=8.4cm,
    axis equal image,
    xlabel={$x$}, ylabel={$y$},
    domain=0:2.6, samples=120,
    xmin=-0.3, xmax=4.2, ymin=-0.3, ymax=4.2,
    axis lines=middle, grid=major, grid style={gray!16},
    legend pos=north west, legend cell align=left,
]
  \addplot[curveblue, very thick, domain=0:2] {x^2};
  \addlegendentry{$f(x)=x^2,\ x\ge0$}
  \addplot[targetgreen!75!black, very thick, domain=0:4] {sqrt(x)};
  \addlegendentry{$f^{-1}(x)=\sqrt{x}$}
  \addplot[gray, dashed, domain=-0.3:4.2] {x};
  \addlegendentry{$y=x$ (mirror)}
\end{axis}
\end{tikzpicture}
\caption{A function and its inverse are mirror images across the line $y=x$. Here
$f(x)=x^2$ (restricted to $x\ge0$ so it is one-to-one) and $f^{-1}(x)=\sqrt x$.
Reflecting across $y=x$ swaps the $x$- and $y$-axes, which is exactly the algebra
of solving $y=x^2$ for $x$.}
\label{ch40:fig-inverse}
\end{figure}

\textbf{Even and odd symmetry.} $f$ is \emph{even} if $f(-x)=f(x)$ (graph
symmetric across the $y$-axis, like $x^2$ or $\cos x$); \emph{odd} if
$f(-x)=-f(x)$ (graph symmetric through the origin, like $x^3$ or $\sin x$). These
symmetries cut work in half: the integral of an odd function over a symmetric
interval $[-a,a]$ is automatically $0$.

\textbf{Monotonicity.} $f$ is \emph{increasing} if $x_1<x_2 \Rightarrow
f(x_1)<f(x_2)$, decreasing if the inequality reverses. A strictly monotonic
function is automatically one-to-one, hence invertible. The derivative's sign
detects monotonicity --- this is the bridge from calculus back to this algebraic
notion.

\begin{trapbox}
$f^{-1}$ means the \emph{inverse function}, not the reciprocal. $\sin^{-1}x$ is the
inverse sine (arcsine), not $1/\sin x$. The reciprocal is written $(\sin x)^{-1}$
or $1/\sin x = \csc x$. The clash of notation is unfortunate but standard; context
and a sharp eye are your defense.
\end{trapbox}

\section{What a Proof Is, and Three Reliable Shapes}
\label{ch40:proof}

A theorem is a true statement; a proof is the evidence that compels belief. But
``evidence'' in mathematics is stricter than anywhere else: a proof is a chain of
steps, each following from the previous by pure logic or from facts already
granted, leaving \emph{no} gap a skeptic could pry open. Checking a million cases
is not a proof if there is a million-and-first. The reward for this strictness is
total certainty: a theorem, once proved, is true forever and everywhere its
hypotheses hold.

\begin{aestheticbox}
Why insist on proof, when experiment and pattern-spotting got us this far?
Because patterns lie. Consider $f(n) = n^2 + n + 41$. For $n=0,1,2,\dots,39$ every
value is prime --- forty cases in a row, overwhelming ``evidence.'' But
$f(40)=40^2+40+41 = 41^2$ is not prime. Forty confirmations could not save a false
conjecture; one proof would have settled it either way. Calculus is built on
infinitely many cases (every $x$, every $\eps$), so we \emph{cannot} check them
one by one. Proof is the only tool that reaches across an infinity at once.
\end{aestheticbox}

We will meet three shapes of proof again and again. Each is a strategy for getting
from hypotheses to conclusion.

\subsection{Direct proof}

The most natural shape: assume the hypothesis, and march forward by valid steps
until you reach the conclusion. Most ``if $P$ then $Q$'' theorems are proved this
way --- start from $P$, deduce, arrive at $Q$.

\begin{example}[A direct proof]
\label{ch40:ex-direct}
\emph{Claim: the sum of two even integers is even.} An even integer is one of the
form $2k$ for some integer $k$ (that is the definition we work from). So let the
two even integers be $m = 2j$ and $n = 2k$. Then
\[
  m + n = 2j + 2k = 2(j+k).
\]
Since $j+k$ is an integer, $m+n$ has the form $2 \times (\text{integer})$, which is
exactly the definition of even. Therefore $m+n$ is even. $\qed$

Notice the rhythm: unfold the definitions of the hypotheses, do honest algebra,
re-fold into the definition of the conclusion. That is the whole craft of a direct
proof.
\end{example}

\subsection{Proof by contradiction}

Sometimes the conclusion is awkward to reach head-on, but its \emph{negation} is
easy to play with. So we assume the opposite of what we want, follow the
consequences, and crash into an impossibility. Since valid reasoning from a true
start can never reach a falsehood, our starting assumption (the negation) must
have been false --- so the original claim is true.

\begin{historybox}
The classic, and one of the most beautiful arguments in all of mathematics, is
Euclid's proof (c.\ 300 BCE) that $\sqrt 2$ is irrational --- that no fraction of
whole numbers equals it. The Pythagoreans are said to have found the existence of
such ``incommensurable'' lengths so disturbing that the discovery was guarded as a
secret; it shattered their creed that all of nature was ratios of whole numbers.
\end{historybox}

\begin{example}[$\sqrt 2$ is irrational --- proof by contradiction]
\label{ch40:ex-sqrt2}
\emph{Suppose, for contradiction, that $\sqrt 2$ is rational.} Then we can write
$\sqrt 2 = p/q$ as a fraction in \emph{lowest terms} --- meaning $p$ and $q$ share
no common factor (any common factor could be cancelled first). Square both sides:
\[
  2 = \frac{p^2}{q^2} \quad\Longrightarrow\quad p^2 = 2q^2.
\]
So $p^2$ is even, which forces $p$ itself to be even (if $p$ were odd, $p^2$ would
be odd). Write $p = 2r$. Substituting,
\[
  (2r)^2 = 2q^2 \;\Longrightarrow\; 4r^2 = 2q^2 \;\Longrightarrow\; q^2 = 2r^2.
\]
Now $q^2$ is even, so $q$ is even too. But then $p$ and $q$ are \emph{both} even
--- they share the common factor $2$ --- contradicting our choice of a fraction in
lowest terms. The contradiction means our supposition was impossible: $\sqrt 2$
cannot be written as such a fraction. Hence $\sqrt 2$ is irrational. $\qed$
\end{example}

\begin{happeningbox}
What really powered that proof? A single squeeze: we assumed a ``smallest'' (lowest
terms) representation and then manufactured a smaller one, which cannot be. This
``no smallest counterexample'' engine --- formally, the well-ordering of the
positive integers, the fact that any nonempty set of them has a least element ---
is the same engine that drives induction, the next technique. Contradiction and
induction are two faces of one principle about $\N$.
\end{happeningbox}

\subsection{Proof by induction}

Induction proves a statement $P(n)$ for \emph{every} natural number $n$ at once,
even though there are infinitely many. The image to hold is a line of dominoes.

\begin{intuitionbox}
Imagine infinitely many dominoes standing in a row, one for each $n=1,2,3,\dots$.
You want to know they all fall. You do not push each one --- there are infinitely
many. Instead you check two things: (1) the \textbf{first} domino falls, and
(2) \textbf{each} domino, if it falls, knocks over the next. Granted those two
facts, every domino must fall: the first falls by (1), which by (2) topples the
second, which topples the third, forever. Induction is exactly this.
\end{intuitionbox}

\begin{focusbox}
\textbf{Principle of mathematical induction.} To prove $P(n)$ holds for all
integers $n \ge 1$, prove two things:
\begin{itemize}
  \item \textbf{Base case:} $P(1)$ is true. (The first domino falls.)
  \item \textbf{Inductive step:} for every $k$, \emph{if} $P(k)$ is true \emph{then}
        $P(k+1)$ is true. (Each domino topples the next.)
\end{itemize}
The assumption ``$P(k)$ is true'' inside the step is called the \emph{inductive
hypothesis}. You are allowed to assume it precisely because you are about to show
it would force the next case.
\end{focusbox}

\begin{example}[Sum of the first $n$ integers]
\label{ch40:ex-induction}
\emph{Claim:} for every $n \ge 1$,
\[
  1 + 2 + 3 + \cdots + n = \frac{n(n+1)}{2}.
\]
\textbf{Base case ($n=1$).} The left side is just $1$; the right side is
$\tfrac{1\cdot 2}{2}=1$. They agree, so $P(1)$ holds.

\textbf{Inductive step.} Assume the formula holds for $n=k$ (the inductive
hypothesis):
\[
  1 + 2 + \cdots + k = \frac{k(k+1)}{2}.
\]
We must show it then holds for $n=k+1$. Add $(k+1)$ to both sides:
\[
  1 + 2 + \cdots + k + (k+1) = \frac{k(k+1)}{2} + (k+1).
\]
Factor $(k+1)$ out of the right-hand side:
\[
  \frac{k(k+1)}{2} + (k+1) = (k+1)\!\left(\frac{k}{2} + 1\right)
  = (k+1)\cdot\frac{k+2}{2} = \frac{(k+1)(k+2)}{2}.
\]
That last expression is exactly the original formula with $n$ replaced by $k+1$.
So $P(k)\Rightarrow P(k+1)$. By induction the formula holds for all $n \ge 1$.
$\qed$
\end{example}

\begin{historybox}
A famous (probably embroidered) story has the young Carl Friedrich Gauss, set the
busywork of adding $1$ through $100$, answering almost instantly: $5050$. His trick
was to pair the sum with itself reversed --- $1+100$, $2+99$, $\dots$, each pair
totalling $101$, and there are $100$ of them, giving $100\cdot 101 = 10100$, so the
sum is half that, $5050$. That pairing is the formula $n(n+1)/2$ seen geometrically,
and it is a second, ``direct,'' proof of the same fact. It is healthy to know a
result two ways: induction certifies it; Gauss's pairing explains \emph{why} the
answer takes that shape.
\end{historybox}

\begin{trapbox}
Induction without a base case proves nothing. ``If $P(k)$ then $P(k+1)$'' alone is
a chain of dominoes that nobody pushed --- the implications can all be valid while
the statement is false for every $n$. Always nail down the base case. Conversely, a
base case with a broken inductive step (a gap where $P(k)$ does not actually force
$P(k+1)$) is equally worthless: one fallen domino, then nothing.
\end{trapbox}

\begin{connectionbox}
Induction is not confined to ``sum'' formulas. It underlies the power rule
$\frac{d}{dx}x^n = nx^{n-1}$ for positive integer $n$ (prove it for $n=1$, then
use the product rule to climb from $x^k$ to $x^{k+1}$), the binomial theorem, and
the convergence arguments behind Taylor series. Whenever a statement is indexed by
a whole number and each case rests on the one before, induction is the tool. It is
the discrete cousin of the ``accumulate one step at a time'' idea that integration
makes continuous.
\end{connectionbox}

\section{The Binomial Theorem: Algebra and Counting Together}
\label{ch40:binomial}

One identity deserves a section of its own because it threads through the entire
book --- it appears in the power rule's derivation, in Taylor's theorem, and in
the definition of $e$. It answers: what is $(a+b)^n$ expanded?

We already know $(a+b)^2 = a^2 + 2ab + b^2$ and $(a+b)^3 = a^3 + 3a^2b + 3ab^2 +
b^3$. The coefficients $1,2,1$ and $1,3,3,1$ are the rows of \emph{Pascal's
triangle}, where each entry is the sum of the two above it:
\[
  \begin{array}{ccccccccc}
   & & & & 1 & & & & \\
   & & & 1 & & 1 & & & \\
   & & 1 & & 2 & & 1 & & \\
   & 1 & & 3 & & 3 & & 1 & \\
   1 & & 4 & & 6 & & 4 & & 1
  \end{array}
\]

\begin{theorem}[Binomial theorem]
\label{ch40:thm-binomial}
For any nonnegative integer $n$,
\[
  (a+b)^n = \sum_{k=0}^{n} \binom{n}{k} a^{n-k} b^{k},
  \qquad \binom{n}{k} = \frac{n!}{k!\,(n-k)!}.
\]
\end{theorem}

\begin{happeningbox}
Where do the coefficients $\binom{n}{k}$ come from? Expanding $(a+b)^n =
(a+b)(a+b)\cdots(a+b)$ ($n$ factors) means choosing, from each factor, either the
$a$ or the $b$, then multiplying the choices. A term $a^{n-k}b^k$ arises every time
we pick $b$ from exactly $k$ of the $n$ factors. The number of ways to choose which
$k$ factors contribute a $b$ is $\binom{n}{k}$ --- ``$n$ choose $k$.'' So the
binomial theorem is really a \emph{counting} statement in algebraic disguise: the
coefficient is the count of ways a given mix of $a$'s and $b$'s can occur. This is
why probability and algebra share the same triangle.
\end{happeningbox}

\begin{example}[Why it matters for the power rule]
\label{ch40:ex-binom-power}
The derivative of $x^n$ from the limit definition needs $(x+h)^n$. The binomial
theorem gives
\[
  (x+h)^n = x^n + n x^{n-1} h + \binom{n}{2} x^{n-2} h^2 + \cdots + h^n.
\]
Subtract $x^n$ and divide by $h$:
\[
  \frac{(x+h)^n - x^n}{h} = n x^{n-1} + \binom{n}{2}x^{n-2} h + \cdots + h^{n-1}.
\]
Every term after the first carries at least one factor of $h$, so as $h \to 0$ they
all vanish, leaving
\[
  \ddx{x^n} = n x^{n-1}.
\]
The binomial theorem isolated the one term that survives the limit. The power rule
is the first coefficient of Pascal's triangle in motion.
\end{example}

\section{Putting It Together: A Worked Gauntlet}
\label{ch40:gauntlet}

To see these tools cooperate, run one expression through several of them.

\begin{example}[Simplify a difference quotient]
\label{ch40:ex-gauntlet}
Simplify $\dfrac{1}{h}\!\left(\dfrac{1}{x+h} - \dfrac{1}{x}\right)$ and find its
limit as $h\to 0$ --- this is the derivative of $1/x$ from scratch.

\textbf{Step 1: combine the inner fractions} over the common denominator
$x(x+h)$:
\[
  \frac{1}{x+h} - \frac{1}{x} = \frac{x - (x+h)}{x(x+h)} = \frac{-h}{x(x+h)}.
\]
\textbf{Step 2: divide by $h$} (multiply by $1/h$); the $h$'s cancel, the move
that makes the limit tractable:
\[
  \frac{1}{h}\cdot\frac{-h}{x(x+h)} = \frac{-1}{x(x+h)}.
\]
\textbf{Step 3: take the limit} $h \to 0$. Now the expression is continuous at
$h=0$, so just substitute:
\[
  \lim_{h\to 0} \frac{-1}{x(x+h)} = \frac{-1}{x \cdot x} = -\frac{1}{x^2}.
\]
So $\ddx{1/x} = -1/x^2$. Every step was pure appendix material --- common
denominators, sign care, cancellation --- yet the result is a genuine calculus
fact. This is the appendix earning its place.
\end{example}

\begin{labbox}
Run the same gauntlet on $\dfrac{\sqrt{x+h}-\sqrt{x}}{h}$. The right tool is the
conjugate: multiply top and bottom by $\sqrt{x+h}+\sqrt{x}$ and use the difference
of squares~\eqref{ch40:diff-squares} to turn the top into $(x+h)-x = h$, which
cancels the $h$ below. You should land on $\dfrac{1}{2\sqrt x}$ in the limit ---
the derivative of $\sqrt x$. Notice it is the same three-act structure: rewrite,
cancel the $h$, substitute.
\end{labbox}

\section*{Exercises}
\addcontentsline{toc}{section}{Exercises}

\begin{exercise}
\label{ch40:exr-factor}
Factor completely over the reals:
(a) $x^3 - 4x$;\quad (b) $x^4 - 16$;\quad (c) $x^2 - 6x + 9$;\quad
(d) $8x^3 - 27$. \\
For each, state how you recognized the pattern (common factor, difference of
squares, perfect square, difference of cubes).
\end{exercise}

\begin{exercise}
\label{ch40:exr-exponents}
Simplify to a single power (no negative or fractional exponents in the final
answer where avoidable):
(a) $\dfrac{a^5 a^{-2}}{a^{-4}}$;\quad
(b) $\bigl(8x^6\bigr)^{1/3}$;\quad
(c) $\dfrac{(2x^2)^3}{4x^4}$;\quad
(d) $\dfrac{x^{1/2}\,x^{1/3}}{x^{-1/6}}$.
Explain at each step which exponent law you are invoking.
\end{exercise}

\begin{exercise}
\label{ch40:exr-quadratic}
Solve $2x^2 - 4x - 3 = 0$ two ways: once with the quadratic formula, and once by
completing the square. Confirm the answers agree, and state the discriminant and
what it tells you about the roots.
\end{exercise}

\begin{exercise}
\label{ch40:exr-inequality}
Solve each inequality, writing the answer in interval notation, and say where (if
anywhere) the ``flip'' rule was needed:
(a) $3 - 2x \le 11$;\quad
(b) $\dfrac{x-1}{x+2} \le 0$;\quad
(c) $x^2 < 5x - 6$.
\end{exercise}

\begin{exercise}
\label{ch40:exr-absval}
Rewrite each as an interval or union of intervals on the number line, and sketch
it:
(a) $\abs{x-3} < 2$;\quad
(b) $\abs{2x+1} \ge 5$;\quad
(c) $\abs{x} = \abs{x-4}$ (interpret this one as ``equidistant from $0$ and $4$''
before computing).
\end{exercise}

\begin{exercise}
\label{ch40:exr-triangle}
Use the triangle inequality to prove the \emph{reverse triangle inequality}
$\bigl|\,\abs{x}-\abs{y}\,\bigr| \le \abs{x-y}$.
(Hint: write $x = (x-y) + y$ and apply $\abs{x+y}\le\abs{x}+\abs{y}$ to get
$\abs{x} \le \abs{x-y} + \abs{y}$; rearrange, then do the same with the roles of
$x$ and $y$ swapped.)
\end{exercise}

\begin{exercise}
\label{ch40:exr-inverse}
Let $f(x) = \dfrac{2x-1}{x+3}$.
(a) Find $f^{-1}$ by solving $y = f(x)$ for $x$.
(b) Verify $f(f^{-1}(y)) = y$.
(c) State the domain of $f$ and the domain of $f^{-1}$, and explain why they are
what they are.
\end{exercise}

\begin{exercise}
\label{ch40:exr-direct}
Give a direct proof that the product of two odd integers is odd. (Start from the
definition: an odd integer has the form $2k+1$.)
\end{exercise}

\begin{exercise}
\label{ch40:exr-contradiction}
Prove by contradiction that $\sqrt 3$ is irrational. Identify the one place where
the argument differs from the $\sqrt 2$ proof of Example~\ref{ch40:ex-sqrt2}
(hint: you will use that ``$3 \mid p^2 \Rightarrow 3 \mid p$,'' which deserves a
sentence of justification).
\end{exercise}

\begin{exercise}
\label{ch40:exr-induction-squares}
Prove by induction that for all $n \ge 1$,
\[
  1^2 + 2^2 + 3^2 + \cdots + n^2 = \frac{n(n+1)(2n+1)}{6}.
\]
Be explicit about your base case, your inductive hypothesis, and the algebra that
turns the $n=k$ formula into the $n=k+1$ formula.
\end{exercise}

\begin{exercise}
\label{ch40:exr-induction-ineq}
Prove by induction that $2^n > n$ for every integer $n \ge 1$. (This is an
\emph{inequality} induction; in the step, add to or multiply the inductive
hypothesis to reach the $k+1$ case.)
\end{exercise}

\begin{exercise}
\label{ch40:exr-binomial}
Use the binomial theorem (Theorem~\ref{ch40:thm-binomial}) to expand $(1+x)^4$,
and verify your coefficients against the corresponding row of Pascal's triangle.
Then use the same expansion to write down the first three terms of the difference
quotient $\bigl((1+h)^4 - 1\bigr)/h$ and read off the derivative of $x^4$ at
$x=1$.
\end{exercise}
